%%%%%%%%%%%%%%%%
%%% exod 19 %%%
%%%%%%%%%%%%%%%%

\Chapter{19}

\Verse{1}
{
	\hP{בַּחֹדֶשׁ הַשְּׁלִישִׁי לְצֵאת בְּנֵי יִשְׂרָאֵל מֵאֶרֶץ מִצְרָיִם בַּיּוֹם הַזֶּה בָּאוּ מִדְבַּר סִינָי׃}
}
{
	\eP{
		In the third month after the exodus of the children of
		Israel from the land of Egypt, on this day, they came to
		the wilderness of Sinai.
	}
}
{
}

\Verse{2}
{
	\hR{וַיִּסְעוּ מֵרְפִידִים וַיָּבֹאוּ מִדְבַּר סִינַי וַיַּחֲנוּ בַּמִּדְבָּר}
	\hE{וַיִּחַן שָׁם יִשְׂרָאֵל נֶגֶד הָהָר׃}
}
{
	\eR{And they traveled from Rephidim and came to the wilderness
		of Sinai and camped in the wilderness.}
	\eE{And Israel camped
		there opposite the mountain.}
}
{
}

\Verse{3}
{
	\hE{וּמֹשֶׁה עָלָה אֶל הָאֱלֹהִים וַיִּקְרָא אֵלָיו יְהֹוָה מִן הָהָר לֵאמֹר כֹּה תֹאמַר לְבֵית יַעֲקֹב וְתַגֵּיד לִבְנֵי יִשְׂרָאֵל׃}
}
{
	\eE{
		And Moses had gone up to God. And YHWH called to him from
		the mountain, saying, “This is what you shall say to the
		house of Jacob and tell to the children of Israel:
	}
}
{
%	\footnote{
%		Moses had gone up to God. While “going up to God” is an intriguing
%		image, the Septuagint text is more probably correct. It reads “up to the
%		mountain of God,” which fits better with the words that follow: “And
%		YHWH called to him from the mountain.”
%	}
}

\Verse{4}
{
	\hE{אַתֶּם רְאִיתֶם אֲשֶׁר עָשִׂיתִי לְמִצְרָיִם וָאֶשָּׂא אֶתְכֶם עַל כַּנְפֵי נְשָׁרִים וָאָבִא אֶתְכֶם אֵלָי׃}
}
{
	\eE{
		‘You’ve seen what I did to Egypt, and I carried you on
		eagles’ wings and brought you to me.
	}
}
{
%	\footnote{
%		on eagles’ wings. The eagle is the highest-flying bird, so there is no
%		threat from above. Eagles fly with their young on their backs.
%	}
}

\Verse{5}
{
	\hE{וְעַתָּה אִם שָׁמוֹעַ תִּשְׁמְעוּ בְּקֹלִי וּשְׁמַרְתֶּם אֶת בְּרִיתִי וִהְיִיתֶם לִי סְגֻלָּה מִכּל הָעַמִּים כִּי לִי כּל הָאָרֶץ׃}
}
{
	\eE{
		And now, if you’ll listen to my voice and observe my
		covenant, then you’ll be a treasure to me out of all the
		peoples, because all the earth is mine.
	}
}
{
%	\footnote{
%		if. This first occurrence of the word segullah (treasure) is not
%		ambiguous in context: The people of Israel will be a treasured
%		possession of God only if they listen and fulfill their covenant. Their
%		status is not based on some intrinsic quality but on their behavior. God
%		is about to make a covenant with them. If they keep their part of it,
%		then they will be segullah. (For more on the matter of being chosen and
%		a treasured people, see the comment on Deut 7:6.) Footnote 19:5. a
%		treasure to me out of all the peoples. Since the verse says explicitly
%		that this treasured status is based on their behavior, and not on their
%		being superior to anyone else, what then is the purpose of their being
%		chosen to observe this covenant? The reason was stated, no less
%		explicitly, to their ancestor Abraham: to be a blessing to all the
%		families of the earth.
%	}
}

\Verse{6}
{
	\hE{וְאַתֶּם תִּהְיוּ לִי מַמְלֶכֶת כֹּהֲנִים וְגוֹי קָדוֹשׁ אֵלֶּה הַדְּבָרִים אֲשֶׁר תְּדַבֵּר אֶל בְּנֵי יִשְׂרָאֵל׃}
}
{
	\eE{
		And you’ll be a kingdom of priests and a holy nation to
		me.’ These are the words that you shall speak to the
		children of Israel.”
	}
}
{
%	\footnote{
%		holy. For a long time we have taught that this word, Hebrew qdš, does
%		not refer to some sacred, elevated condition but rather that it denotes
%		being separate, set aside for some particular purpose. But, really, it
%		does in fact refer to a special consecrated state. When God tells Moses
%		at the bush, “Take off your shoes from your feet, because the place on
%		which you’re standing: it’s holy ground,” this means more than that the
%		ground is “set aside.” It means “holy” in just the way that people
%		naturally understand it on first reading. The ground has a special
%		quality deriving from the realm of the divine which makes shoes
%		inappropriate there. And in the present context, “holy nation” does not
%		mean that Israel will be separate from other nations. It means that
%		Israel will be consecrated if the people will live the life that their
%		divine covenant requires of them. (See also the comment on Exod 29:37.)
%	}
}

\Verse{7}
{
	\hE{וַיָּבֹא מֹשֶׁה וַיִּקְרָא לְזִקְנֵי הָעָם וַיָּשֶׂם לִפְנֵיהֶם אֵת כּל הַדְּבָרִים הָאֵלֶּה אֲשֶׁר צִוָּהוּ יְהֹוָה׃}
}
{
	\eE{
		And Moses came and called the people’s elders and set
		before them all these words that YHWH had commanded him.
	}
}
{
}

\Verse{8}
{
	\hE{וַיַּעֲנוּ כל הָעָם יַחְדָּו וַיֹּאמְרוּ כֹּל אֲשֶׁר דִּבֶּר יְהֹוָה נַעֲשֶׂה וַיָּשֶׁב מֹשֶׁה אֶת דִּבְרֵי הָעָם אֶל יְהֹוָה׃}
}
{
	\eE{
		And all the people responded together, and they said,
		“We’ll do everything that YHWH has spoken.” And Moses
		brought back the people’s words to YHWH.
	}
}
{
}

\Verse{9}
{
	\hE{וַיֹּאמֶר יְהֹוָה אֶל מֹשֶׁה הִנֵּה אָנֹכִי בָּא אֵלֶיךָ בְּעַב הֶעָנָן בַּעֲבוּר יִשְׁמַע הָעָם בְּדַבְּרִי עִמָּךְ וְגַם בְּךָ יַאֲמִינוּ לְעוֹלָם וַיַּגֵּד מֹשֶׁה אֶת דִּבְרֵי הָעָם אֶל יְהֹוָה׃}
}
{
	\eE{
		And YHWH said to Moses, “Here, I am coming to you in a
		mass of cloud for the purpose that the people will hear
		when I am speaking with you, and they will believe in you
		as well forever.” And Moses told the people’s words to
		YHWH.
	}
}
{
}

\Verse{10}
{
	\hJ{וַיֹּאמֶר יְהֹוָה אֶל מֹשֶׁה לֵךְ אֶל הָעָם וְקִדַּשְׁתָּם הַיּוֹם וּמָחָר וְכִבְּסוּ שִׂמְלֹתָם׃}
}
{
	\eJ{
		And YHWH said to Moses, “Go to the people and consecrate
		them today and tomorrow; and they shall wash their clothes
	}
}
{
}

\Verse{11}
{
	\hJ{וְהָיוּ נְכֹנִים לַיּוֹם הַשְּׁלִישִׁי כִּי בַּיּוֹם הַשְּׁלִשִׁי יֵרֵד יְהֹוָה לְעֵינֵי כל הָעָם עַל הַר סִינָי׃}
}
{
	\eJ{
		and be ready for the third day, because on the third day
		YHWH will come down on Mount Sinai before the eyes of all
		the people.
	}
}
{
}

\Verse{12}
{
	\hJ{וְהִגְבַּלְתָּ אֶת הָעָם סָבִיב לֵאמֹר הִשָּׁמְרוּ לָכֶם עֲלוֹת בָּהָר וּנְגֹעַ בְּקָצֵהוּ כּל הַנֹּגֵעַ בָּהָר מוֹת יוּמָת׃}
}
{
	\eJ{
		And you shall limit the people all around, saying, ‘Watch
		yourselves about going up in the mountain and touching its
		edge. Anyone who touches the mountain shall be put to
		death.
	}
}
{
}

\Verse{13}
{
	\hJ{לֹא תִגַּע בּוֹ יָד כִּי סָקוֹל יִסָּקֵל אוֹ יָרֹה יִיָּרֶה אִם בְּהֵמָה אִם אִישׁ לֹא יִחְיֶה בִּמְשֹׁךְ הַיֹּבֵל הֵמָּה יַעֲלוּ בָהָר׃}
}
{
	\eJ{
		A hand shall not touch him, but he shall be stoned or
		shot. Whether animal or man, he will not live.’ At the
		blowing of the horn they shall go up the mountain.”
	}
}
{
}

\Verse{14}
{
	\hJ{וַיֵּרֶד מֹשֶׁה מִן הָהָר אֶל הָעָם וַיְקַדֵּשׁ אֶת הָעָם וַיְכַבְּסוּ שִׂמְלֹתָם׃}
}
{
	\eJ{
		And Moses went down from the mountain to the people. And
		he consecrated the people, and they washed their clothes.
	}
}
{
}

\Verse{15}
{
	\hJ{וַיֹּאמֶר אֶל הָעָם הֱיוּ נְכֹנִים לִשְׁלֹשֶׁת יָמִים אַל תִּגְּשׁוּ אֶל אִשָּׁה׃}
}
{
	\eJ{
		And he said to the people, “Be ready for three days. Don’t
		come close to a woman.”
	}
}
{
%	\footnote{
%		he said to the people, “Don’t come close to a woman.” If he is speaking
%		to “the people,” which includes both men and women, why does he say,
%		“Don’t come close to a woman”? Like many passages, this may have two
%		opposite interpretations. It may be a case of male chauvinism in
%		language, in which an address to “the people” means the men. Or it may
%		reflect a perception that a command to abstain from sex for three days
%		needs to be particularly directed to men because men are more likely
%		than women to violate the instruction.
%	}
}

\Verse{16}
{
	\hJ{וַיְהִי בַיּוֹם הַשְּׁלִישִׁי בִּהְיֹת הַבֹּקֶר}
	\hE{וַיְהִי קֹלֹת וּבְרָקִים וְעָנָן כָּבֵד עַל הָהָר וְקֹל שֹׁפָר חָזָק מְאֹד וַיֶּחֱרַד כּל הָעָם אֲשֶׁר בַּמַּחֲנֶה׃}
}
{
	\eJ{And it was on the third day, when it was morning,}
	\eE{and it
		was: thunders and lightning and a heavy cloud on the
		mountain, and a sound of a horn, very strong. And the
		entire people that was in the camp trembled.}
}
{
%	\footnote{
%		a heavy cloud. The word “heavy” recurs, but this time it refers to
%		something divine.
%	}
}

\Verse{17}
{
	\hE{וַיּוֹצֵא מֹשֶׁה אֶת הָעָם לִקְרַאת הָאֱלֹהִים מִן הַמַּחֲנֶה וַיִּתְיַצְּבוּ בְּתַחְתִּית הָהָר׃}
}
{
	\eE{
		And Moses brought out the people toward God from the camp,
		and they stood up at the bottom of the mountain.
	}
}
{
}

\Verse{18}
{
	\hJ{וְהַר סִינַי עָשַׁן כֻּלּוֹ מִפְּנֵי אֲשֶׁר יָרַד עָלָיו יְהֹוָה בָּאֵשׁ וַיַּעַל עֲשָׁנוֹ כְּעֶשֶׁן הַכִּבְשָׁן וַיֶּחֱרַד כּל הָהָר מְאֹד׃}
}
{
	\eJ{
		And Mount Sinai was all smoke because YHWH came down on it
		in fire, and its smoke went up like the smoke of a
		furnace, and the whole mountain trembled greatly.
	}
}
{
}

\Verse{19}
{
	\hE{וַיְהִי קוֹל הַשֹּׁפָר הוֹלֵךְ וְחָזֵק מְאֹד מֹשֶׁה יְדַבֵּר וְהָאֱלֹהִים יַעֲנֶנּוּ בְקוֹל׃}
}
{
	\eE{
		And the sound of the horn was getting much stronger. Moses
		would speak, and God would answer him in a voice.
	}
}
{
%	\footnote{
%		in a voice. Here, Hebrew bqôl can mean “in a voice” or “with sound”
%		(i.e., aloud) or “with thunder.” In v. 16 the word refers first to
%		thunder sounds and then to the sound of a horn. But the divine statement
%		in v. 9 that “the people will hear when I am speaking with you” suggests
%		that the people will hear spoken words between God and Moses, and that
%		is what now happens in v. 19.
%	}
}

\Verse{20}
{
	\hJ{וַיֵּרֶד יְהֹוָה עַל הַר סִינַי אֶל רֹאשׁ הָהָר וַיִּקְרָא יְהֹוָה לְמֹשֶׁה אֶל רֹאשׁ הָהָר וַיַּעַל מֹשֶׁה׃}
}
{
	\eJ{
		And YHWH came down on Mount Sinai, at the top of the
		mountain, and YHWH called to Moses at the top of the
		mountain, and Moses went up.
	}
}
{
}

\Verse{21}
{
	\hJ{וַיֹּאמֶר יְהֹוָה אֶל מֹשֶׁה רֵד הָעֵד בָּעָם פֶּן יֶהֶרְסוּ אֶל יְהֹוָה לִרְאוֹת וְנָפַל מִמֶּנּוּ רָב׃}
}
{
	\eJ{
		And YHWH said to Moses, “Go down. Warn the people in case
		they break through to YHWH, to see, and many of them fall.
	}
}
{
}

\Verse{22}
{
	\hJ{וְגַם הַכֹּהֲנִים הַנִּגָּשִׁים אֶל יְהֹוָה יִתְקַדָּשׁוּ פֶּן יִפְרֹץ בָּהֶם יְהֹוָה׃}
}
{
	\eJ{
		And also let the priests who approach YHWH consecrate
		themselves, or else YHWH will break out against them.”
	}
}
{
}

\Verse{23}
{
	\hJ{וַיֹּאמֶר מֹשֶׁה אֶל יְהֹוָה לֹא יוּכַל הָעָם לַעֲלֹת אֶל הַר סִינָי כִּי אַתָּה הַעֵדֹתָה בָּנוּ לֵאמֹר הַגְבֵּל אֶת הָהָר וְקִדַּשְׁתּוֹ׃}
}
{
	\eJ{
		And Moses said to YHWH, “The people is not able to go up
		to Mount Sinai, because you warned us, saying, ‘Limit the
		mountain and consecrate it.’”
	}
}
{
}

\Verse{24}
{
	\hJ{וַיֹּאמֶר אֵלָיו יְהֹוָה לֶךְ רֵד וְעָלִיתָ אַתָּה וְאַהֲרֹן עִמָּךְ וְהַכֹּהֲנִים וְהָעָם אַל יֶהֶרְסוּ לַעֲלֹת אֶל יְהֹוָה פֶּן יִפְרץ בָּם׃}
}
{
	\eJ{
		And YHWH said to him, “Go. Go down. Then you’ll come up,
		you and Aaron with you. And let the priests and the people
		not break through to come up to YHWH, or else He’ll break
		out against them.”
	}
}
{
}

\Verse{25}
{
	\hJ{וַיֵּרֶד מֹשֶׁה אֶל הָעָם וַיֹּאמֶר אֲלֵהֶם׃}
}
{
	\eJ{And Moses went down to the people, and he said it to them.}
}
{
}
